\paragraph{A permutation-trace refinement.}
The forced central fermion $G$ in the pseudo-unitary Galois conjugate of
$R8\_3\_941$ supplies another necessary test.  Since $N_G$ is an
involutive permutation matrix and has monodromy eigenvalues $\pm1$, its
number of fixed simples equals $P-N$, where $P$ and $N$ count the two
charges.  Only negative-charge simples can be fixed by balancing.
The ten known distinguished objects and translates form five nonfixed
pairs.  If $p_j$ counts the positive-charge pairs and $q_j$ the remaining
negative-charge high simples of restriction type $j$, then
\[
 P=8+2\sum_jp_j,\qquad N=2+\sum_jq_j,
 \qquad
 \sum_j(q_j\bmod2)\le P-N\le\sum_jq_j,
 \quad P-N\equiv\sum_jq_j\pmod2.
\]
The complete $93$-allocation domain retains $71$ allocations.  Intersecting
with the sixteen exact paired-phase survivors removes allocations $4$
and $86$: they respectively require $-2$ fixed objects, or thirteen
fixed objects among only seven eligible ones.

Products of fusion matrices also have nonnegative integral traces.  The
receipt tests all $420$ quantities
$\operatorname{Tr}(N_G^e\prod_{i=1}^4N_{A_i}^{a_i})$, with
$e\in\{0,1\}$ and $\sum_i a_i\le6$, on every allocation, using the
full simultaneous modular eigenvalues.  An actual $Z(\mathrm{Ising})$
integer-matrix control checks the normalization.  No further survivor is
removed.  The final intersection leaves fourteen rational paired-cone
witnesses, eleven of them with exact nonnegative integer witnesses, and
three unresolved integer-feasibility questions.  These are finite-system
survivals, not constructions of fusion categories.
